\cotefile{78}
\scanpage{103}
b) Tout sommet $s=\varphi_S(\underline{s})$ appartient à
exactement deux arêtes. Soient en effet
$\underline a,\underline a'\in\underline A$ les arêtes incidentes à $\underline s$, $a,a'$
leurs images dans $A$, \deleted{\unsure{on}} sait que $a\ne a'$,
$s\in a$, $s\in a'$. Montrons que si $a_0\in\varphi_A(\underline a_0)$
\deleted{est} contient $s$, alors $a_0=a$ ou $a_0=a'$,
\deleted{\unsure{Prenons}} Soient $\deleted{\underline s_0}\underline{\unsure{t|s}}_0,\underline t_1$ les sommets incidents
à $\underline a_0$, \deleted{donc} $t_0,t_1\in S$ leurs images, donc
\deleted{\unsure{$s,t,t'$ sont|$s_0\ne s_1$, et}} \inserted{$t_0\ne t_1$, et} $a_0=\text{\unsure{dr|Dr}}\,(t_0,t_1)\ni s$, et par a) on
doit avoir \deleted{\unsure{$t\in s$}} $s=t_0$ ou $s=t_1$, \deleted{\unsure{et par|disons}}
\deleted{$s$ est inci} $s=\underline t_0$. Soit $\underline r=(\underline t_0,\underline a_0)$,
on a $\underline s=\underline u^{n}\underline t_0$, pour $n\in\mathbb Z$ convenable,
donc $u^nt_0=\deleted{s}=s_0$ \deleted{\unsure{d'où}} $u^n=\mathrm{id}$, donc
\deleted{\illegible{1}} $u^n(a_0)\deleted{\unsure{=ua}} \inserted{=a_0}$. Or $\underline{s}_0$ étant incident
à $\underline a_0$, $\underline s=\underline u^n(\underline s_0)$ l'est à $\underline u^n(\underline a_0)$, donc
$\underline u^n(\underline a_0)$ \deleted{\unsure{$\underline a$}} est égal à $\underline a$ ou $\underline a'$, donc
$u^n(a_0)$ \inserted{$=a_0$} est égal à $a$ ou $a'$, \deleted{donc} cqfd.

c) \deleted{\unsure{les}} $S$ engendre $E$ : connu.

d) $G=\mathrm{Aut}(E,S,A)$ est transitif sur
l'en\inserted{$R$} des repères de $(S,A)$. Dans
la dém. de a) \inserted{(ou comme cas immédiat de a))} on a vu que $R$ est
% archive sheet number: 102
